\documentclass[a4paper,11pt]{article}
\usepackage[utf8]{inputenc}
\usepackage[T1]{fontenc}
\usepackage[french]{babel}
\usepackage{geometry}
\geometry{left=2.5cm,right=2.5cm,top=2.5cm,bottom=2.5cm}
\usepackage{titlesec}
\usepackage{enumitem}
\usepackage{hyperref}

\titleformat{\section}{\large\bfseries}{\thesection}{1em}{}

\begin{document}

\begin{center}
    {\huge \textbf{Projet de Recherche}} \\
    \bigskip
    {\large \textbf{Ismail AISSAOUI}} \\
    Candidature au Doctorat : Programme Engineering Digital Twin (EDT) \\
    Sujet : Composition et Configuration d'Architectures de Jumeaux Numériques
\end{center}

\bigskip

\section{Introduction et Motivation}
La création de jumeaux numériques performants nécessite l'assemblage de composants logiciels et matériels hétérogènes. Cependant, la configuration de ces architectures pour répondre à des exigences contradictoires (consommation d'énergie, latence, précision, coût) reste un défi majeur. Ce projet de recherche vise à développer des méthodes formelles et des outils d'optimisation pour automatiser la composition d'architectures de jumeaux numériques.

\section{Objectifs de Recherche}
Le projet s'articulera autour de trois axes principaux :
\begin{enumerate}[label=\textbf{Axe \arabic*:}]
    \item \textbf{Modélisation de l'Espace de Conception (DSE) :} Définir un méta-modèle permettant de décrire les variantes architecturales d'un jumeau numérique et leurs impacts sur les exigences non-fonctionnelles.
    \item \textbf{Optimisation Multi-Objectif via Métaheuristiques :} Utiliser des algorithmes tels que l'optimisation par essaim particulaire (PSO) ou les algorithmes génétiques pour naviguer dans l'espace des configurations et identifier le front de Pareto optimal.
    \item \textbf{Validation sur Cas d'Usage Industriels :} Appliquer les méthodes développées à des scénarios réels (ex: maintenance prédictive, gestion de l'énergie) pour valider la robustesse des architectures générées.
\end{enumerate}

\section{Méthodologie}
Nous adopterons une approche itérative combinant l'ingénierie dirigée par les modèles (MDE) et l'apprentissage par renforcement ou l'optimisation métaheuristique. L'objectif est de fournir aux concepteurs un outil d'aide à la décision capable de suggérer la meilleure configuration matérielle/logicielle en fonction de contraintes changeantes en temps réel.

\section{Conclusion}
Ce travail de thèse contribuera à la construction de jumeaux numériques plus durables et plus adaptables, renforçant ainsi la position de la recherche française dans le domaine de l'industrie 4.0 et du programme national EDT.

\end{document}
